\documentclass[12pt]{article}
\usepackage[T1]{fontenc}
\usepackage[utf8]{inputenc}
\usepackage{lmodern}
\usepackage{textcomp}
\usepackage{enumitem}
\usepackage{geometry}
\geometry{margin=1in}
\begin{document}
\begin{center}
Answer Key - Constitutional Law - Undergraduate Level Assessment 7 \
\small (Answers are ordered exactly as in the question paper.)
\end{center}

\section*{Multiple Choice}
\begin{enumerate}[label=\arabic*.]
\item \textbf{Answer:} B
\\ \textit{Options:} A) The coastal State; B) The flag State; C) All States enjoy such jurisdiction; D) The International Tribunal for the Law of the Sea
\\ \textit{Note:} The flag State exercises jurisdiction over crimes committed on board vessels because international law grants the state whose flag the vessel flies the authority to regulate activities and enforce laws on that ship, regardless of its location on the high seas.
\item \textbf{Answer:} B
\\ \textit{Options:} A) An act is jure imperii when undertaken by an international organisation; B) An act is jure imperii when undertaken in an official State capacity; C) All acts undertaken by State officials are acts jure imperii; D) An act is jure imperii when undertaken by a State corporation
\\ \textit{Note:} An act jure imperii refers to actions taken by a state in its sovereign capacity, distinguishing them from private acts, thereby establishing the legal principle that such acts are performed within the framework of official governmental authority.
\item \textbf{Answer:} B
\\ \textit{Options:} A) Hugo Grotius; B) Hobbes; C) Locke; D) Rousseau
\\ \textit{Note:} The correct answer is Hobbes because he argued in his social contract theory that natural law arises from human nature and the necessity of societal living, emphasizing the inherent need for order and governance to avoid chaos in human interactions.
\item \textbf{Answer:} A
\\ \textit{Options:} A) The law is too complex to be codified.; B) Codification ossifies the law.; C) Bentham fails to take account of the criminal law.; D) Significant portions of the law are already codified.
\\ \textit{Note:} The statement that "the law is too complex to be codified" effectively refutes Bentham's argument because it highlights the inherent intricacies and nuances of legal systems that may resist simplification into a single code, suggesting that codification could lead to oversimplification and misinterpretation of the law.
\item \textbf{Answer:} A
\\ \textit{Options:} A) Natural Law; B) Analytical; C) Historical; D) Sociological
\\ \textit{Note:} The Natural Law School of jurisprudence asserts that law is derived from inherent moral principles and universal truths about what is "correct," emphasizing that laws must align with ethical standards.
\end{enumerate}

\section*{True / False}
\begin{enumerate}[label=\arabic*.]
\setcounter{enumi}{5}
\item \textbf{Answer:} True
\\ \textit{Options:} A) True; B) False
\\ \textit{Note:} Hobbes argues in his social contract theory that in the absence of a governing authority to maintain peace, individuals revert to a state of nature where they possess unrestricted rights to pursue their self-interests, including the right to take the lives of others.
\item \textbf{Answer:} False
\\ \textit{Options:} A) True; B) False
\\ \textit{Note:} The European Convention on Human Rights primarily emphasizes civil and political rights, such as the right to life, freedom from torture, and the right to a fair trial, rather than focusing on economic and social rights.
\item \textbf{Answer:} True
\\ \textit{Options:} A) True; B) False
\\ \textit{Note:} Austin is considered a 'naive empiricist' because he emphasizes the practical application of laws based on observable experiences rather than engaging with abstract conceptual frameworks, which aligns with a pragmatic approach to understanding legal principles.
\item \textbf{Answer:} False
\\ \textit{Options:} A) True; B) False
\\ \textit{Note:} The statement is false because, in the original position as conceived by John Rawls, individuals would prioritize a fair distribution of social goods that balances both liberty and equality to ensure justice as fairness.
\item \textbf{Answer:} True
\\ \textit{Options:} A) True; B) False
\\ \textit{Note:} Sovereignty encompasses the authority of a state to govern itself and make decisions independently, free from external control or influence.
\end{enumerate}

\section*{Long Form Response}
\begin{enumerate}[label=\arabic*.]
\setcounter{enumi}{10}
\item \textbf{Answer:} Territorial jurisdiction is a fundamental principle in constitutional law that establishes the geographic boundaries within which a legal authority can exercise its power. This concept is significant as it ensures that laws and governance apply consistently within a defined area, thereby promoting order and predictability in legal proceedings. The implications of territorial jurisdiction are profound, as they delineate the limits of governmental authority, influence the enforcement of laws, and affect the rights of individuals and entities within those boundaries. Furthermore, this principle underpins the relationship between different jurisdictions, particularly in cases involving cross- border legal issues, thereby highlighting the need for cooperation and conflict resolution among various legal systems. Ultimately, understanding territorial jurisdiction is essential for comprehending how laws are applied and enforced within a constitutional framework.
\\ \textit{Note:} correct because it accurately defines territorial jurisdiction as the geographic boundaries for legal authority, emphasizing its importance in ensuring consistent application of laws, establishing governmental limits, and affecting law enforcement.
\item \textbf{Answer:} The entities entitled to request advisory opinions from the International Court of Justice (ICJ) include the United Nations General Assembly and the Security Council, which can seek guidance on any legal question. Additionally, other UN organs and specialized agencies may request advisory opinions if authorized by the General Assembly, specifically concerning legal questions within their operational scope. These requests play a significant role in shaping international law by providing authoritative interpretations that can influence state behavior and international governance. The advisory opinions serve not only to clarify legal ambiguities but also to promote adherence to international legal standards, thereby strengthening the rule of law globally.
\\ \textit{Note:} correct because it accurately identifies the entities entitled to request advisory opinions from the ICJ, emphasizing the role of the UN General Assembly, Security Council, and authorized UN organs in shaping international law through these legal inquiries.
\end{enumerate}

\vfill
\noindent\textit{End of Answer Key}
\end{document}